\chapter{La Danza Final: Interacción Fluido-Estructura}
\label{chap:fsi}

%==================================================================================
\section{Cuando el Fluido y el Sólido Bailan Juntos}
%==================================================================================

Hasta este punto, hemos acoplado la física \textit{dentro} de un mismo dominio: el calor afectaba a la mecánica en un sólido, y la velocidad afectaba al calor en un fluido. Ahora damos el paso final y más complejo: acoplar dos físicas a través de la \textbf{frontera móvil} que las separa. Bienvenidos al mundo de la \textbf{Interacción Fluido-Estructura (FSI)}.

En un problema FSI, el flujo de un fluido ejerce fuerzas sobre una estructura, causando que esta se deforme. Esta deformación, a su vez, altera la forma del dominio del fluido, lo que cambia el patrón de flujo y las fuerzas que ejerce. Este ciclo de retroalimentación es el corazón de algunos de los fenómenos de ingeniería más complejos y fascinantes.

Resolver problemas FSI es uno de los pináculos de la simulación multifísica, ya que requiere un solver robusto tanto para la dinámica de fluidos como para la mecánica de sólidos, y una estrategia inteligente para manejar el acoplamiento y el movimiento de la malla en la interfaz.

\begin{center}
\begin{tabular}{|p{0.9\textwidth}|}
\hline
\textbf{CONTEXTO INGENIERIL} \\
La FSI es fundamental en el diseño de sistemas donde las estructuras flexibles interactúan con flujos:
\begin{itemize}
    \item \textbf{Ingeniería Aeroespacial:} El fenómeno del \textit{flutter} en las alas de los aviones, una vibración autoexcitada que puede llevar a fallos catastróficos.
    \item \textbf{Ingeniería Biomédica:} El flujo de sangre a través de las válvulas del corazón o en arterias flexibles; el flujo de aire en los pulmones.
    \item \textbf{Ingeniería Civil:} El efecto del viento sobre puentes colgantes (como el infame caso del puente de Tacoma Narrows), edificios altos y torres.
    \item \textbf{Energía:} El flujo de agua a través de las palas de una turbina hidroeléctrica.
\end{itemize}
\\
\hline
\end{tabular}
\end{center}

%==================================================================================
\section{La Física del Acoplamiento en una Interfaz Móvil}
\label{sec:fsi-theory}
%==================================================================================
Un problema FSI consta de tres componentes: el dominio del fluido, el dominio del sólido y la interfaz que los une.

%----------------------------------------------------------------------------------
\subsection{Las Ecuaciones en Cada Dominio}
%----------------------------------------------------------------------------------
\begin{itemize}
    \item \textbf{En el Dominio del Fluido (\(\Omega_f\)):} El flujo se rige por las ecuaciones de \textbf{Navier-Stokes} (Ecuaciones~\eqref{eq:ns-momentum} y \eqref{eq:ns-continuity}).
    \item \textbf{En el Dominio del Sólido (\(\Omega_s\)):} La deformación se rige por la \textbf{ecuación de la elasticidad} (o elastodinámica si se consideran los efectos inerciales de la estructura), Ecuación~\eqref{eq:equilibrium}.
\end{itemize}

%----------------------------------------------------------------------------------
\subsection{Las Condiciones de Acoplamiento en la Interfaz}
%----------------------------------------------------------------------------------
La magia de la FSI ocurre en la interfaz \(\Gamma_{FSI}\) que separa el fluido y el sólido. En esta frontera, dos condiciones deben cumplirse simultáneamente:
\begin{enumerate}
    \item \textbf{Acoplamiento Cinemático (Condición de No Deslizamiento):} La velocidad del fluido en la interfaz debe ser igual a la velocidad de la estructura en esa misma interfaz. El fluido "se pega" a la estructura.
    \begin{equation}
        \bm{u}_f = \frac{\partial \bm{d}_s}{\partial t} \quad \text{en } \Gamma_{FSI}
        \label{eq:fsi-kinematic}
    \end{equation}
    donde \(\bm{u}_f\) es la velocidad del fluido y \(\bm{d}_s\) es el desplazamiento del sólido.
    \item \textbf{Acoplamiento Dinámico (Equilibrio de Tracciones):} Las fuerzas que el fluido ejerce sobre la estructura deben ser iguales y opuestas a las fuerzas que la estructura ejerce sobre el fluido (tercera ley de Newton).
    \begin{equation}
        \bm{\sigma}_f \cdot \bm{n}_f = \bm{\sigma}_s \cdot \bm{n}_s \quad \text{en } \Gamma_{FSI}
        \label{eq:fsi-dynamic}
    \end{equation}
    donde \(\bm{\sigma}_f\) y \(\bm{\sigma}_s\) son los tensores de esfuerzo del fluido y del sólido, respectivamente.
\end{enumerate}

%----------------------------------------------------------------------------------
\subsection{El Desafío del Dominio Móvil: La Formulación ALE}
%----------------------------------------------------------------------------------
El mayor desafío numérico de la FSI es que el dominio del fluido \(\Omega_f\) cambia de forma a medida que la estructura se deforma. Resolver las ecuaciones de Navier-Stokes en un dominio que se mueve y deforma constantemente es complejo.

La técnica más común para abordar esto es la formulación \textbf{Arbitraria Lagrangeana-Euleriana (ALE)}. La idea es la siguiente:
\begin{itemize}
    \item En una descripción \textbf{Lagrangeana} (típica de los sólidos), los nodos de la malla se mueven junto con el material.
    \item En una descripción \textbf{Euleriana} (típica de los fluidos), la malla está fija y el fluido se mueve a través de ella.
    \item La descripción \textbf{ALE} es un híbrido. Los nodos de la malla del fluido pueden moverse, pero no necesariamente con la misma velocidad que el fluido.
\end{itemize}
En la práctica, hacemos que los nodos de la malla en la interfaz \(\Gamma_{FSI}\) se muevan con la estructura (para satisfacer la condición cinemática), y este movimiento se propaga suavemente hacia el interior de la malla del fluido, a menudo resolviendo una ecuación de elasticidad artificial para la propia malla. Esto requiere modificar las ecuaciones de Navier-Stokes para incluir la velocidad de la malla, \(\bm{w}\), un tema avanzado que se trata en detalle en textos como \cite{Zienkiewicz2014}.

%==================================================================================
\section{Estrategias de Solución para Problemas FSI}
\label{sec:fsi-strategy}
%==================================================================================
Existen dos grandes familias de estrategias para resolver el sistema FSI acoplado:

\begin{itemize}
    \item \textbf{Enfoques Monolíticos:} Se construye una única y gigantesca matriz de sistema que incluye todas las incógnitas (velocidad y presión del fluido, y desplazamiento del sólido) y se resuelve todo de una vez en cada paso de tiempo. Son muy precisos y estables, pero extremadamente exigentes computacionalmente y difíciles de implementar.
    \item \textbf{Enfoques Particionados (o Segregados):} Se utiliza una estrategia de "divide y vencerás", muy similar a la que usamos para la convección forzada. En cada paso de tiempo:
    \begin{enumerate}
        \item Se resuelven las ecuaciones del fluido, asumiendo una posición de la estructura.
        \item Se calculan las fuerzas del fluido sobre la estructura.
        \item Se pasan estas fuerzas al solver de la estructura.
        \item Se resuelven las ecuaciones del sólido para encontrar su nueva deformación.
        \item Se actualiza la posición de la malla del fluido según la deformación de la estructura.
        \item Se repite el proceso hasta la convergencia.
    \end{enumerate}
    Estos métodos son mucho más flexibles y fáciles de implementar, ya que permiten reutilizar solvers de fluidos y sólidos existentes, y son el enfoque más común en el software comercial.
\end{itemize}
Con estas herramientas conceptuales, estamos preparados para abordar la simulación de sistemas donde la danza entre el fluido y la estructura define la física del problema.